% 星际行星（综述）
% license CCBYSA3
% type Wiki

本文根据 CC-BY-SA 协议转载翻译自维基百科\href{https://en.wikipedia.org/wiki/Rogue_planet}{相关文章}。

流浪行星（rogue planet），亦称自由漂浮行星（free-floating planet，FFP）或孤立的行星质量天体（isolated planetary-mass object，iPMO），是指一种具有行星质量、但在引力上不与任何恒星或棕矮星发生束缚关系的星际天体$^{[1][2][3][4]}$。

流浪行星可能起源于其形成所在的行星系统，并在随后被抛射出去；它们也可能在行星系统之外独立形成。仅银河系中就可能存在数十亿到数万亿颗流浪行星，这一数量范围有望由即将投入使用的南希·格雷斯·罗曼空间望远镜进一步加以精确限定$^{[5][6]}$。流浪行星进入太阳系的概率极低，更不用说对地球生命构成直接威胁——在未来一千年内发生这种情况的概率约为一万亿分之一$^{[7]}$。

一些行星质量天体可能以类似恒星的方式形成，国际天文学联合会已提议将此类天体称为亚棕矮星（sub-brown dwarfs）$^{[8]}$。一个可能的例子是 Cha~110913$-$773444，它可能是在被抛射后成为一颗流浪行星，也可能是独立形成并成为一颗亚棕矮星$^{[9]}$。
\subsection{术语}
最早的两篇发现论文分别使用了“孤立的行星质量天体”（iPMO）$^{[10]}$和“自由漂浮行星”（FFP）$^{[11]}$这两个名称。大多数天文学论文使用其中之一$^{[12][13][14]}$。“流浪行星”这一术语更常用于引力微透镜研究，而这类研究也经常使用 FFP 这一术语$^{[15][16]}$。面向公众的新闻稿则可能采用不同的名称。例如，2021 年发现至少 70 颗 FFP 时，不同新闻稿分别使用了“流浪行星”$^{[17]}$、“无恒星行星”$^{[18]}$、“游荡行星”$^{[19]}$以及“自由漂浮行星”$^{[20]}$等称呼。
\subsection{发现}
孤立的行星质量天体（isolated planetary-mass objects，iPMO）最早于 2000 年由英国研究团队 Lucas \& Roche 使用英国红外望远镜（UKIRT）在猎户座星云中发现$^{[11]}$。同年，西班牙研究团队 Zapatero Osorio 等人利用凯克望远镜的光谱观测，在 $\sigma$~猎户座星团中发现了 iPMO$^{[10]}$。对猎户座星云中这些天体的光谱研究于 2001 年发表$^{[21]}$。目前，这两个欧洲研究团队均因其准同步发现而受到认可$^{[22]}$。1999 年，日本研究团队 Oasa 等人在变色龙座 I 云中发现了一些天体$^{[23]}$，这些天体随后在 2004 年由美国研究团队 Luhman 等人通过光谱观测得到了确认$^{[24]}$。
\subsection{观测}
发现自由漂浮行星主要有两种技术：直接成像和引力微透镜。
\subsubsection{引力微透镜}
日本大阪大学的天体物理学家角广隆宏（Takahiro Sumi）及其同事，隶属于天体物理学引力微透镜观测计划（Microlensing Observations in Astrophysics）和光学引力透镜实验（Optical Gravitational Lensing Experiment）合作组，于 2011 年发表了他们关于引力微透镜的研究成果。他们使用位于新西兰约翰山天文台的 1.8 米（5 英尺 11 英寸）MOA-II 望远镜，以及位于智利拉斯坎帕纳斯天文台的华沙大学 1.3 米（4 英尺 3 英寸）望远镜，对银河系中约 5000 万颗恒星进行了观测。他们共发现了 474 次微透镜事件，其中有 10 次事件的持续时间足够短，表明它们可能是质量约为木星质量、且在其附近没有伴随恒星的行星。研究人员根据这些观测结果估计，在银河系中，每一颗恒星大约对应近两颗木星质量的流浪行星$^{[25][26][27]}$。另有一项研究提出了一个大得多的数量估计，认为银河系中的流浪行星数量可能是恒星数量的多达 100{,}000 倍，不过该研究所涵盖的假想天体质量远小于木星$^{[28]}$。2017 年，华沙大学天文台的 Przemek~Mróz 及其同事在统计样本数量为 2011 年研究六倍的基础上开展的一项研究表明，银河系中每一颗主序星所对应的木星质量自由漂浮行星或宽轨道行星的上限为 0.25 颗$^{[29]}$。

2020 年 9 月，天文学家利用引力微透镜技术首次报告探测到一颗地球质量的流浪行星（命名为 OGLE-2016-BLG-1928），该天体不与任何恒星发生引力束缚，并在银河系中自由漂浮$^{[16][30][31]}$。
\subsubsection{直接成像}


